\documentclass[10pt,a4paper]{amsart}
\usepackage[margin=19mm]{geometry}
\usepackage{amsmath,amssymb,mathtools}
\usepackage[scr=rsfs,cal=euler]{mathalfa}
\usepackage{xfrac}
\usepackage[unicode,hidelinks,bookmarks=false]{hyperref}
\newcommand{\ie}{i.e.\ }
\newcommand{\cf}{cf.\ }
\newcommand{\resp}{respectively\ }
\newcommand{\Subsection}{\subsection}
\providecommand{\nobreakdash}{\nobreak\hbox{-}}
\setlength{\emergencystretch}{2em}
\multlinegap0pt
\usepackage{bm}
\usepackage{bbm}
\usepackage{dsfont}
\usepackage{xspace}
\usepackage{cancel}
\usepackage{mathtools}
\usepackage{amsthm}
\makeatletter
\@ifundefined{thm}{\newtheorem{thm}{Thm}}{}
\@ifundefined{lemma}{\newtheorem{lemma}[thm]{Lemma}}{}
\makeatother
\DeclareMathOperator{\SO}{\textnormal{SO}} % constants
\DeclareMathOperator{\Vb}{\mathbf{V}} % bold vector
\DeclareMathOperator{\Wb}{\mathbf{W}} % bold vector
\DeclareMathOperator{\Xb}{\mathbf{X}} % bold vector
\DeclareMathOperator{\cP}{\mathcal{P}} %
\DeclareMathOperator{\cS}{\mathcal{S}} %
\DeclareMathOperator{\one}{\mathbf{1}} % indicator function
\DeclareMathOperator{\sphere}{\mathbb{S}} % sphere
\DeclareMathOperator{\xb}{\mathbf{x}} % bold vector
\def\bY{\overline{Y}}
\def\rr{\mathbb{R}}
\title[Unimodality for Radon partitions of random vectors]{Unimodality for Radon partitions of random vectors}
\author{Swee Hong Chan, Gil Kalai, Bhargav Narayanan, Natalya Ter-Saakov, Moshe White}
\date{Source: arXiv:2507.01353v1 (CC BY 4.0)}
\begin{document}
\maketitle

\noindent\emph{Source and licence.} Unimodality for Radon partitions of random vectors. Version arXiv:2507.01353v1 (CC BY 4.0), distributed under the Creative Commons Attribution 4.0 International licence, as recorded in the arXiv licence metadata for this exact version and shown on the abstract page https://arxiv.org/abs/2507.01353v1. Full rights notice: licenses/arxiv-2507.01353-CC-BY-4.0.txt.\par\smallskip

\noindent\textit{Editorial scope.} Counted unit: the Lemma whose source label is \texttt{lem:iso} in the article (source0.tex lines 541--543, source label \texttt{lem:iso}), delivered together with the article's own complete original proof of it (source0.tex lines 545--573, one \texttt{proof} environment of 29 lines). No mathematics is written, repaired or re-derived: statement and proof are the article's own text line for line. Required context retained uncounted: none. Every definition, display and label of those lines is retained unchanged. Omitted from this package and not redistributed: 1--540, 544--544, 574--882. Nothing inside a retained excerpt is omitted or altered: every retained body is delivered exactly as the article's lines are written, so no omission receipt of any kind accompanies this package. Citations: the retained excerpts carry no citation command at all, so no bibliography entry of the article is transcribed and no reference is redistributed. Reference adaptation: none was needed and none was made; every reference inside the delivered unit resolves inside it, so no unresolved reference or citation is produced. Printed slips kept literal: no printed word, symbol, hypothesis or number of the counted statement or of its proof was altered. Layout and class: the article's own class is \texttt{amsart} with packages including geometry, amsmath, amssymb, amsthm, bbm, mathtools, comment, amsfonts, enumitem, hyperref, euscript, tikz, babel, adjustbox; the wrapper is the lane's pinned \texttt{amsart} editorial wrapper at 10pt on A4, which restates the theorem-like environments the retained text needs and adds the article's own macro definitions that the retained text needs (\texttt{\textbackslash SO}, \texttt{\textbackslash Vb}, \texttt{\textbackslash Wb}, \texttt{\textbackslash Xb}, \texttt{\textbackslash bY}, \texttt{\textbackslash cP}, \texttt{\textbackslash cS}, \texttt{\textbackslash one}, \texttt{\textbackslash rr}, \texttt{\textbackslash sphere}, \texttt{\textbackslash xb}), with extra wrapper packages (bm, bbm, dsfont, xspace, cancel, mathtools). The delivered numbering is the unit's own. Assets: no retained span contains a figure environment, a graphics-inclusion command, a PDF-page inclusion, a table or a TikZ picture; no image file is copied into this package, the package declares no asset dependency and no image file is redistributed with it. Licence: version arXiv:2507.01353v1 of this article is distributed under the Creative Commons Attribution 4.0 International licence, as recorded in the arXiv licence metadata for this exact version and shown on the abstract page https://arxiv.org/abs/2507.01353v1. A full rights notice is delivered at licenses/arxiv-2507.01353-CC-BY-4.0.txt. Characters: every retained excerpt is pure ASCII, so the delivered engine, XeLaTeX with its default Latin Modern text font, reports no missing character.
\begin{lemma}\label{lem:iso}
	The random variables $\cS_n$ 	and $\cP_n$ are equal in distribution.
\end{lemma}

\begin{proof}
	Let $K$ be the subset of the unit sphere $\sphere^{n-1}$ defined by
	\[ K = \left \{ \xb \in \sphere^{n-1} : \one \cdot \xb=0 \right \}, \]
	where $\one=(1,\dots, 1) \in \rr^n$, and let $\mu$ be the uniform probability measure on $K$. Note that this measure is invariant under the action of the group $G \leq \SO(n)$ given by
	\[ G = \left \{ A \in \SO(n):  \one \cdot A \in \{\one,-\one\} \right\}, \]
    where $\SO(n)$ is the special orthogonal group.
	It is easily seen that the action of $G$ on $K$ is transitive, so $\mu$ is the unique probability measure on $K$ that is invariant under the action of $G$.

	We now express the law of $\cS_n$ in terms of $\mu$. Let $\Vb=(V_1,\dots, V_{n}) \in K$ be a vector that satisfies
	\[
		V_1 \Xb_1 + \cdots + V_{n} \Xb_{n} = 0.
	\]
	Note that distinct $i,j \in [n]$ belong to the same part of the Radon partition of $\Xb_1,\dots, \Xb_{n}$ if and only if $V_i$ and $V_j$ have the same sign. Also note that $\Vb$ is unique up to a sign change almost surely since $\Xb_1,\dots, \Xb_{n}$ are in general position almost surely. From the two vectors $\Vb$ and $-\Vb$, we choose one uniformly at random, and we denote this random vector by $\Wb=(W_1,\dots, W_{n})$. It then follows from the construction that
	\begin{equation}\label{eq:iso-1}
		\cS_n = \{ i \in [n] \mid W_i \geq 0 \}.
	\end{equation}
	Now, note that the measure induced by $(\Xb_1,\dots, \Xb_{n})$ is invariant under the action of $\SO(n)$ as $\Xb_1,\dots, \Xb_n$ are independent standard normal random vectors. It then follows that the measure induced by $\Wb$ is also invariant under the action of $G$. Therefore, we conclude that $\Wb$ is distributed according to $\mu$.


	Next, we express the law of $\cP_n$ in terms of $\mu$. Let $\Vb'=(V_1',\dots, V_n') \in \rr^n$ be the random vector given by
	\[ V_i' = \bY - Y_i \]
	for $i \in [n]$,  and let $\Wb'=(W_1',\dots, W_n')$ be the vector in $K$ given by
	\[ \Wb' = \frac{1}{\sqrt{\Vb' \cdot \Vb'}} \Vb'. \]
	It then follows from the construction that
	\begin{equation}\label{eq:iso-2}
		\cP_n = \{ i \in [n] :  V_i' \geq 0 \} = \{ i \in [n] :  W_i' \geq 0 \}.
	\end{equation}
	Since the measure induced by $(Y_1,\dots, Y_n)$ is invariant under the action of $\SO(n)$, it follows that the measure induced by $\Wb'$ is unique under the action of $G$. We conclude that $\Wb'$ is distributed according to $\mu$ as well, and the lemma now follows from comparing~\eqref{eq:iso-1} and~\eqref{eq:iso-2}.
\end{proof}

\end{document}
